\documentclass[11pt,a4paper]{article}
\usepackage[T1]{fontenc}
\usepackage[utf8]{inputenc}
\usepackage[osf,sc]{mathpazo}
\usepackage[scaled=0.94]{helvet}
\usepackage[margin=18mm,top=16mm,bottom=17mm]{geometry}
\usepackage{microtype}
\usepackage{textcomp}
\usepackage{xcolor}
\usepackage{enumitem}
\usepackage{titlesec}
\usepackage{fancyhdr}
\usepackage{lastpage}
\usepackage{hyperref}

\definecolor{ink}{HTML}{202A34}
\definecolor{muted}{HTML}{566473}
\definecolor{accent}{HTML}{1A5276}
\definecolor{rulecolor}{HTML}{D8E0E6}
\hypersetup{colorlinks=true,urlcolor=accent,linkcolor=accent,
  pdftitle={Fran Bartolic | Curriculum Vitae},pdfauthor={Fran Bartolic}}
\color{ink}
\setlength{\parindent}{0pt}
\setlength{\parskip}{0pt}
\setlength{\emergencystretch}{2em}
\linespread{1.04}
\setlist[itemize]{leftmargin=1.25em,label=\textcolor{muted}{\textbullet},
  itemsep=1pt,topsep=3pt,parsep=0pt,partopsep=0pt}
\titleformat{\section}{\large\sffamily\bfseries\color{accent}}{}{0pt}{}[\vspace{2pt}\color{rulecolor}\titlerule]
\titlespacing*{\section}{0pt}{11pt}{6pt}
\pagestyle{fancy}
\fancyhf{}
\renewcommand{\headrulewidth}{0pt}
\fancyfoot[L]{\footnotesize\sffamily\color{muted}Fran Bartoli\'{c}\enspace /\enspace Curriculum vitae}
\fancyfoot[R]{\footnotesize\sffamily\color{muted}\thepage\ /\ \pageref*{LastPage}}
\setlength{\footskip}{20pt}

% Shared entry formatting keeps content separate from layout.
\newcommand{\cvitem}[2]{%
  \par\noindent\begin{minipage}[t]{24mm}
    % Enter horizontal mode before setting color to preserve baseline alignment.
    \raggedright\small\sffamily\leavevmode\color{muted}#1
  \end{minipage}\hfill\begin{minipage}[t]{\dimexpr\linewidth-28mm\relax}
    \raggedright #2
  \end{minipage}\par\vspace{4pt}}
\newcommand{\cventry}[6]{%
  \cvitem{#1}{\textbf{#2}\enspace|\enspace #3\\[-1pt]
    {\small\sffamily\color{muted}#4}\if\relax\detokenize{#5}\relax\else\enspace #5\fi
    \par\vspace{2pt}#6}}

\begin{document}
{\fontsize{30}{34}\selectfont Fran Bartoli\'{c}}\par
\vspace{5pt}
{\large\sffamily Research Scientist\enspace|\enspace Google DeepMind}\par
\vspace{3pt}
{\sffamily\color{muted}AI weather forecasting\enspace\textperiodcentered\enspace New York, USA}\par
\vspace{7pt}
{\small\sffamily
franbartolic at google.com\enspace\textperiodcentered\enspace
\href{https://fbartolic.github.io}{fbartolic.github.io}\enspace\textperiodcentered\enspace
\href{https://github.com/fbartolic}{GitHub}\enspace\textperiodcentered\enspace
\href{https://www.linkedin.com/in/fbartolic/}{LinkedIn}}\par
\vspace{3pt}

\section{Experience}
\cventry{Jul 2026--\\Present}{Research Scientist}{Google DeepMind}
{New York, USA}{}{Developing AI weather forecasting models for WeatherNext.}
\cventry{Dec 2022--\\Jul 2026}{Research Scientist}{Salient Predictions}
{Remote (Croatia), then New York, USA}{}{
    \begin{itemize}
        \item Developed probabilistic weather forecasting models across medium-range, subseasonal, and seasonal timescales.
        \item Developed methods for combining forecasts from multiple models, including neural-network approaches.
        \item Co-developed \href{https://arxiv.org/abs/2601.03753}{GEM-2} and \href{https://arxiv.org/abs/2608.06241}{GEM-3}, efficient global weather models for probabilistic forecasting.
    \end{itemize}}
\cventry{Jun--Dec 2021}{Research Resident}{Cervest}
{London, UK}{}{
    \begin{itemize}
        \item Built probabilistic models of the urban heat island effect using geospatial data.
    \end{itemize}}
\cventry{Feb--Jun 2020}{Research Analyst}{Center for Computational Astrophysics, Flatiron Institute}
{New York, USA}{}{
    \begin{itemize}
        \item Developed a probabilistic \href{https://speakerdeck.com/fbartolic/inferring-a-time-dependent-map-of-io-from-occultations-and-phase-curves}{model} to reconstruct time-varying planetary surface maps from light curves using nonnegative matrix factorization and variational inference.
    \end{itemize}}

\section{Education}
\cventry{2017--2023}{Ph.D. in Astrophysics}{University of St Andrews}
{St Andrews, Scotland}{}{
    \begin{itemize}
        \item Developed Bayesian methods for gravitational microlensing and planetary surface mapping from astronomical time series.
        \item Created \href{https://github.com/fbartolic/caustics}{caustics}, an open-source JAX package for differentiable computation of binary- and triple-lens microlensing light curves.
    \end{itemize}}
\cventry{2015--2017}{M.Sc. in Physics with Astrophysics}{University of Rijeka}
{Rijeka, Croatia}{}{
    \begin{itemize}
        \item Conducted \href{https://github.com/fbartolic/master_thesis}{thesis research} on circumbinary exoplanet dynamics at Lund Observatory, Sweden.
    \end{itemize}}
\cventry{2012--2015}{B.Sc. in Physics}{University of Split}
{Split, Croatia}{}{
    \begin{itemize}
        \item Studied whether tidal engulfment explains the scarcity of close-in giant planets around evolved stars.
        \item Exchange semester at Lund University, Sweden.
    \end{itemize}}

\newpage
\section{Publications}
\cvitem{2026}{S. Levang, \textbf{F. Bartoli\'{c}}, T. Dickinson, C. Dwelle, P. Rauba, V. Cikojevi\'{c}. \emph{Timestep-Conditioned Transformers for Global Weather Forecasting}, \href{https://arxiv.org/abs/2608.06241}{arXiv:2608.06241}.}
\cvitem{2026}{P. Rauba, V. Cikojevi\'{c}, \textbf{F. Bartoli\'{c}}, S. Levang, T. Dickinson, C. Dwelle. \emph{Probabilistic Transformers for Joint Modeling of Global Weather Dynamics and Decision-Centric Variables}, \href{https://arxiv.org/abs/2601.03753}{arXiv:2601.03753}.}
\cvitem{2022}{\textbf{F. Bartoli\'{c}}, R. Luger, D. Foreman-Mackey, R. R. Howell, J. A. Rathbun. \emph{Occultation mapping of Io\textquotesingle s surface in the near-infrared I: Inferring static maps}, The Planetary Science Journal. doi:10.3847/PSJ/ac2a3e.}
\cvitem{2021}{R. Luger, E. Agol, \textbf{F. Bartoli\'{c}}, D. Foreman-Mackey. \emph{Analytic Light Curves in Reflected Light: Phase Curves, Occultations, and Non-Lambertian Scattering for Spherical Planets and Moons}, \href{https://arxiv.org/abs/2103.06275}{arXiv:2103.06275}.}
\cvitem{2020}{N. Golovich, W. Dawson, \textbf{F. Bartoli\'{c}}, et al. \emph{A Reanalysis of Public OGLE-III and IV Gravitational Microlensing Events}, \href{https://arxiv.org/abs/2009.07927}{arXiv:2009.07927}.}
\cvitem{2018}{V. Bozza, E. Bachelet, \textbf{F. Bartoli\'{c}}, T. M. Heintz,
A. R. Hoag, and M. Hundertmark.  \emph{VBBINARYLENSING: a public package for microlensing light-curve  computation.}, 2018, MNRAS, 479, 5157. doi:10.1093/mnras/sty1791}

\section{Patents}
\cvitem{2026}{S. Levang, \textbf{F. Bartoli\'{c}}. \emph{Multi-model blending of probabilistic weather forecasts}, \href{https://patents.justia.com/patent/20260003100}{US Patent Application 2026/0003100 A1} (filed 06/2025).}
\cvitem{2026}{S. Levang, \textbf{F. Bartoli\'{c}}. \emph{Multi-model blending via a neural network for probabilistic weather forecasts}, \href{https://patents.justia.com/patent/20260003101}{US Patent Application 2026/0003101 A1} (filed 06/2025).}
\cvitem{2026}{S. Ridge, V. Cikojevi\'{c}, \textbf{F. Bartoli\'{c}}. \emph{Long-term weather forecasting framework using statistical and machine learning modeling}, \href{https://patents.justia.com/patent/20260104533}{US Patent Application 2026/0104533 A1} (filed 10/2025).}

\section{Awards}
\cvitem{2019}{\emph{Arthur Maitland Prize} for the best talk, University of St Andrews.}
\cvitem{2015}{\emph{Dean's Award for undergraduate academic excellence},
University of Split.}

\section{Additional information}
\cvitem{Languages}{English (fluent), Croatian (native).}
\cvitem{Interests}{Cooking, reading, AI, history.}
\end{document}
